\documentclass[11pt]{article}
\usepackage[T1]{fontenc}
\usepackage[utf8]{inputenc}
\usepackage{lmodern,amsmath,amssymb,microtype,needspace}
\usepackage[margin=1in]{geometry}
\usepackage[colorlinks=true,urlcolor=blue]{hyperref}
\setlength{\emergencystretch}{3em}
\setlength{\parindent}{0pt}
\setlength{\parskip}{.3em}
\title{Response to reviewers\\MPS-D-26-01262}
\author{John T. Foster}
\date{29 October 2026}
\begin{document}
\maketitle
We thank the reviewers for their careful assessment. The revised manuscript is
\emph{Apparent non-associated plasticity from distinct isochoric and volumetric
mechanisms}. The responses below address every substantive point in the
provided reports from Reviewers 1 and 3. Reviewer wording is reproduced with
line-break hyphenation and math typesetting normalized. All section, equation, figure, table, and page references
refer to the revised clean manuscript, not the marked-up version.

The main revisions do three things. They explain why the independent model
remains associated and how the coupled model produces a different flow
direction. They identify the elastic energies that the split allows. They add
an exact finite-deformation example that a reader can check by hand.

\Needspace{6\baselineskip}
\section{R1-General}
\textbf{Comment.}\begin{quote}\small
I find this manuscript quite well written and I like the fact that it is sharp and quite clearly focused. However, I am not sure that the obtained result (basically, what the Author writes at the end of page 6 and about Eq. (23), its extension to a specific finite-deformation framework, and its link with Drucker-Prager plasticity through a kinematic constraint under specific conditions) is enough to deserve publication in JMPS. Perhaps, the Author could convince me by adding a thorough discussion on relevant implications of the proposed way of seeing non-associated plasticity as the combination of two independent associated plasticity mechanisms. Also, more complex options than the Flory split might help in appreciating this proposal.
\end{quote}
\textbf{Response.} The revised discussion explains what separating the plastic
rates means for a confined specimen that dilates while it shears. Shear
supplies the work, and that work pays for two things: expansion against the
confining pressure and dissipation. The dilation slope sets the expansion work,
while the difference between the friction and dilation slopes sets the
dissipation. We connect this work budget explicitly to Collins and Houlsby's
stress-dependent dissipation construction and explain which finite-deformation
factor carries each rate.

We also separate two requirements that are often treated as one: nonnegative
dissipation, and classical maximum plastic dissipation over a fixed convex
elastic domain. The independent model satisfies both with respect to its
complete domain. In the coupled model, normality holds in generalized-force
space when the pressure is held fixed. Because the dissipative resistance
depends on pressure, the normal in the physical stresses differs. Two
independent associated mechanisms and one coupled shear-dilation mechanism
therefore represent different material responses.

The exact homogeneous finite-deformation example states the elastic energy, the
preload, the loading history, and the resulting stress path, volume change,
dissipation, and small-strain limit. The energy-class derivation identifies the
extra state variable, conjugate force, and evolution law that an independently
energetic porosity model would need. We explain the physical role of that
extension and use the Flory closure for the worked example.

\textbf{Revised locations.} Section~1 (p.~1); Section~3.1 (p.~5); Section~5.1 (p.~10); Section~6.1 (p.~12); Section~7 (p.~14); Section~8 (p.~15).
\Needspace{6\baselineskip}
\section{R1-A}
\textbf{Comment.}\begin{quote}\small
The Author writes that the multiplicative decomposition of the deformation gradient is "classical". However, it does not necessarily represent the real behaviour of the materials here concerned with (while it should hold in crystal plasticity, as demonstrated by Reina and Conti, 2014, JMPS). Because of this, for instance, I disagree with the terminology "true-plastic" velocity used for \(\bar L^p\); moreover, about this terminology, I would, in case, use it for the difference between the velocity gradient and \(\bar L^e=\dot{\bar F}^e(\bar F^e)^{-1}\), which is not what here is defined as \(\bar L^p\).
\end{quote}
\textbf{Response.} We agree that an established multiplicative decomposition
does not by itself show what the microscopic mechanisms are in these materials.
The introduction now calls it a modeling choice and cites Reina and Conti
(2014). Their crystalline-slip construction does not cover the granular or
pore-collapse case, so the manuscript treats the present factors as
constitutive assumptions that require material-specific evidence.

We now call the mechanism ``isochoric plastic deformation'' and describe
\(\bar{\mathbf L}^p\) as its velocity gradient in the intermediate
configuration. We keep ``true deformation'' for the kinematic factor in
Drumheller's split. The plastic rate reaches the current configuration as
\(\bar{\mathbf F}^e\bar{\mathbf L}^p(\bar{\mathbf F}^e)^{-1}\). The text now
explains why this is not the difference between the total spatial velocity
gradient and the elastic true-deformation rate: that difference also contains
the distension contribution. The added reference is Reina and Conti,
\emph{JMPS} 67 (2014), 40--61, DOI:
\href{https://doi.org/10.1016/j.jmps.2014.01.014}{10.1016/j.jmps.2014.01.014}.

\textbf{Revised locations.} Section~1 (p.~1); Section~2 (p.~3).
\Needspace{6\baselineskip}
\section{R1-B}
\textbf{Comment.}\begin{quote}\small
It is not clear from the text that the first equation in section 2.1 straightforwardly follows from Eqs. (1)--(4). Up to that point, it is just standard finite-deformation kinematics, such that I find sentences like "The present single-phase theory does not carry volume fraction or mass exchange, so an application must instead prescribe a constitutive evolution or constraint for a." obscure, on the one hand, and obvious, on the other hand. In fact, at the end of the day, the Author just considers the trivial assumption in Eq. (5). Analogously, Eq. (8) comes straightforwardly from finite-deformation elasto-plasticity relying on the multiplicative decomposition of the deformation gradient. So, up to the end of section 3 I find it difficult to appreciate either novelties or a sufficiently clear and thorough discussion of the literature.
\end{quote}
\textbf{Response.} The closure section now identifies the determinant-rate
relation as Jacobi's identity and states what a closure has to supply. We
connect the scalar distension to Herrmann's porous-compaction description and
to Drumheller's mixture kinematics, and we distinguish the absolute
bulk-to-solid volume ratio from a reference-normalized change. When volume
fraction is retained, mass balance identifies the distension and a constitutive
law determines its evolution.

Sections~2 and~3 establish the kinematics and the thermodynamic restrictions
that the model uses. The rescaling calculation shows that, under the stated
power balance, the two unrestricted elastic factors enter the stored energy
only through their combined elastic deformation. The Flory specialization
follows from constrained isochoric variations. The revised discussion states
the contribution of the paper directly: a physical reading of the plastic-power
terms, the activation and energetic restrictions they imply, and a worked
finite-deformation realization. It also compares the construction with Collins
and Houlsby's formulation and with classical maximum plastic dissipation.

\textbf{Revised locations.} Section~2.1 (p.~4); Section~3.1 (p.~5); Section~5.1 (p.~10); Section~7 (p.~14).
\Needspace{6\baselineskip}
\section{R1-C-i}
\textbf{Comment.}\begin{quote}\small
The Author uses "associative" in the title only, while all along the manuscript He uses "associated". Does it make sense?
\end{quote}
\textbf{Response.} The title now uses ``non-associated,'' matching the
terminology used throughout the manuscript.

\textbf{Revised locations.} Title (p.~1).
\Needspace{6\baselineskip}
\section{R1-C-ii}
\textbf{Comment.}\begin{quote}\small
Eq. (1) does not characterise at all the two-point tensors A and \(\bar F\) and I do not understand the sentence below it.
\end{quote}
\textbf{Response.} We now introduce the split through a physical distinction:
compression of the solid material, and change in the volume assigned to it. The
paragraph following Eq.~(1) defines the two maps in order. The true-deformation
tensor takes reference line elements into an intermediate configuration that
represents deformation of the solid material. The distension tensor takes them
from that configuration into the current spatial tangent space. We state that
these are local two-point maps whose global compatibility is a separate
condition, and we point to the closure that identifies the factors.

The construction is explicitly grounded in Drumheller (2000), Sections 3--4 and
8.2. We also cite Herrmann's (1969) pressure--distension model to explain the
scalar bulk-to-solid volume ratio, and we distinguish that scalar
interpretation from the full tensorial map. The closure section explains the
reference normalization and the additional evolution information required when
volume fraction is retained.

\textbf{Revised locations.} Eq.~1 (p.~3); Section~2 (p.~3); Section~2.1 (p.~4).
\Needspace{6\baselineskip}
\section{R1-C-iii}
\textbf{Comment.}\begin{quote}\small
In the definitions in the equations in between Eq. (2) and Eq. (3) there are symbols that, although are clear, are left undefined.
\end{quote}
\textbf{Response.} We now define the identity tensor, the elastic and plastic
superscripts, the logarithmic distensions, and the material time derivative.
The superscripts are defined where the factors are split, and the remaining
definitions follow immediately after the isotropic-factor display.

\textbf{Revised locations.} Section~2 (p.~3); Eq.~3 (p.~3).
\Needspace{6\baselineskip}
\section{R1-C-iv}
\textbf{Comment.}\begin{quote}\small
The first time the Kirchhoff stress is used (at about the beginning of section 3) is an undefined symbol.
\end{quote}
\textbf{Response.} At first use we now define the Kirchhoff stress as
\(\boldsymbol\tau=\mathbf P\mathbf F^T=J\boldsymbol\sigma\), and we name the
first Piola and Cauchy stresses before the frame-indifference discussion.

\textbf{Revised locations.} Section~3 (p.~5).
\Needspace{6\baselineskip}
\section{R1-C-v}
\textbf{Comment.}\begin{quote}\small
In several quotations in the text the year is missing.
\end{quote}
\textbf{Response.} Years have been added to the narrative author attributions,
including Drucker, Bland, Drucker--Prager, Collins--Houlsby, Drumheller,
Foster and coauthors, Srinivasa, and Ulloa and coauthors. The full reference
entries retain publication years. Numeric citation callouts are retained, as
required by the manuscript's bibliography style.

\textbf{Revised locations.} Section~1 (p.~1); Section~2 (p.~3); Section~5.1 (p.~10).
\Needspace{6\baselineskip}
\section{R1-C-vi}
\textbf{Comment.}\begin{quote}\small
I do not understand the reasoning about the first equation on page 5: doesn't the expression between round brackets vanish anyway?
\end{quote}
\textbf{Response.} Yes. If the unrestricted tensor stress restriction is
imposed first, the whole bracket vanishes and its deviatoric projection adds
nothing. The revised text states that the Flory calculation is instead an
alternative derivation through a constrained process. There, only
volume-preserving elastic variations are admissible, so an extension of the
elastic energy away from the isochoric manifold need satisfy the projected
condition and not the stronger unrestricted tensor equation. The scalar
volumetric condition supplies the spherical response.

\textbf{Revised locations.} Section~3.1 (p.~5); Eq.~11 (p.~6).
\Needspace{6\baselineskip}
\section{R1-C-vii}
\textbf{Comment.}\begin{quote}\small
First line of section 4: "Their" seems to me incorrect, as Eq. (9) refers to only one of the two mechanisms.
\end{quote}
\textbf{Response.} The opening of Section~4 now says ``For the isochoric
mechanism, the equivalent stress ...,'' rather than assigning both equivalent
quantities to both mechanisms.

\textbf{Revised locations.} Section~4 (p.~7); Eq.~13 (p.~7).
\Needspace{6\baselineskip}
\section{R1-C-viii}
\textbf{Comment.}\begin{quote}\small
Kuhn-Tucker -> Karush-Kuhn-Tucker.
\end{quote}
\textbf{Response.} ``Kuhn--Tucker'' has been replaced by ``Karush--Kuhn--Tucker.''

\textbf{Revised locations.} Section~4 (p.~7); Eq.~16 (p.~7).
\Needspace{6\baselineskip}
\section{R3-1}
\textbf{Comment.}\begin{quote}\small
The interpretation of association deserves a more explicit distinction between the independent and coupled constructions. In Section 4, the sum of active associated flow directions remains within the normal cone of the complete elastic domain, even if it is not normal to either individual surface. In Section 5, the generalized-force derivative is instead taken with pressure held fixed, before identifying the volumetric force with minus the physical pressure. Please distinguish these two settings and explain how they support the term “apparent non-associativity.” In particular, clarify that the proportional volumetric-to-shear plastic-rate relation is an imposed constitutive constraint. A sketch of the yield domain and flow directions would help communicate this useful interpretation.
\end{quote}
\textbf{Response.} The independent model is now described consistently as
associated with its complete convex elastic domain. At a fixed hardening state,
the combined rate is a nonnegative sum of the outward normals to the active
yield surfaces, so it belongs to the normal cone of the full domain.

In the coupled model, the derivative with respect to the generalized forces is
taken at fixed pressure. Replacing the volumetric force by minus the physical
pressure then gives a yield surface whose derivative also contains the pressure
dependence of the dissipative strength. We now identify the dilation ratio
explicitly as an imposed constitutive constraint, and the loading consistency
condition includes the change of pressure. The discussion relates this
distinction to maximum plastic dissipation and explains the shear work needed
to drive dilation under confinement.

We explain the geometry with the explicit force-space equations and the loading
comparison rather than with a schematic sketch. These specify which variables
are held fixed and show the difference in physical response directly.

\textbf{Revised locations.} Section~4 (p.~7); Section~5 (p.~8); Eq.~23 (p.~9); Section~7 (p.~14); Section~8 (p.~15).
\Needspace{6\baselineskip}
\section{R3-2}
\textbf{Comment.}\begin{quote}\small
Please clarify the admissible free-energy class in Section 3. When the elastic true-deformation and distension rates are treated as independent, Eqs. (7a,b) jointly imply invariance of the energy under a compensating rescaling of these elastic variables. How does this restriction accommodate the interpretation of distension as an independently evolving microstructural quantity? A short derivation and a clear distinction between the unrestricted formulation and the Flory-closed specialization would strengthen the thermodynamic presentation.
\end{quote}
\textbf{Response.} We agree, and the revised text now derives the energy class
explicitly. Contracting the unrestricted tensor restriction with the elastic
true-deformation factor and combining it with the scalar restriction gives
\(\psi_{,\bar{\mathbf F}^e}:\bar{\mathbf F}^e=3\psi_{,\zeta^e}\). Integrating
this identity along the compensating rescaling gives
\(\psi=\widehat\psi(e^{\zeta^e/3}\bar{\mathbf F}^e,\kappa_d,\kappa_v)\).

Under the stated power balance, changing the two unrestricted factors in
compensating ways therefore does not change the energy, and the factors do not
store independent distension energy. We have removed the implication that a
closure or evolution law alone can create that energetic degree of freedom. An
independently energetic porosity or other microstructure requires an additional
state variable and conjugate force together with an appropriate balance or
evolution law, and that force is not generally just minus the pressure. The
Flory-closed specialization is distinguished by its constrained isochoric
process directions and its projected stress condition.

\textbf{Revised locations.} Section~3.1 (p.~5); Eq.~8 (p.~5); Eq.~9 (p.~6); Eq.~10 (p.~6); Eq.~11 (p.~6).
\Needspace{6\baselineskip}
\section{R3-3}
\textbf{Comment.}\begin{quote}\small
The small-strain proportional-loading calculation is informative. To substantiate the finite-deformation extension in Section 6, please add one worked homogeneous material-point example within the stated compressive loading regime. Specify the elastic free energy, initial state, parameters and loading history, and report the stress path, plastic volume change and plastic dissipation. It would also be useful to verify preservation of the isochoric plastic determinant and recovery of the small-strain limit. This example would make the formulation easier to assess and implement without requiring a large-scale boundary-value simulation.
\end{quote}
\textbf{Response.} The new example in Section~6.1 uses an objective isotropic
Hencky-type elastic energy, a stress-free initial state, an elastic
hydrostatic/shear preload, and a prescribed proportional coaxial
logarithmic-strain history. All elastic constants and plastic parameters are
stated. The diagonal factors commute, so the solution is closed-form and
Table~1 is evaluated from it directly.

The example reports Kirchhoff and Cauchy stress relations, total and plastic
volume ratios, the logarithmic plastic volume change, the instantaneous
dissipation, and the accumulated dissipated energy. The accumulated value is
\(46.6941\,\mathrm{kJ/m^3}\) of reference volume at the final state. The
example verifies \(\det\bar{\mathbf F}^p=1\) identically and gives the exact
power balance. Scaling the preload and the accumulated logarithmic strains
recovers the small-strain stress-rate law and multiplier, and the Cauchy and
Kirchhoff measures are compared explicitly at three scales. Together these
provide analytical checks for a material-point implementation along the stated
loading history.

\textbf{Revised locations.} Section~6.1 (p.~12); Eq.~31 (p.~12); Eq.~34 (p.~13); Eq.~35 (p.~13); Table~1 (p.~14); Eq.~36 (p.~13); Eq.~37 (p.~13).
\Needspace{6\baselineskip}
\section{R3-4}
\textbf{Comment.}\begin{quote}\small
Please distinguish Kirchhoff and Cauchy pressure invariants explicitly, for example by using different subscripts, and state whether the reported dissipation is measured per unit reference or current volume.
\end{quote}
\textbf{Response.} The stress definition now relates the Kirchhoff and Cauchy
stresses explicitly. The finite-deformation notation uses \(p\) for the
Kirchhoff pressure and \(p/J\) for the corresponding Cauchy pressure. The
worked example reports both pressures and gives the shear-invariant conversion
for its coaxial path. The thermodynamic section now states that dissipation is
power per reference volume, and division by \(J\) gives the value per current
volume. The example and its table state their units and volume basis.

\textbf{Revised locations.} Section~3 (p.~5); Section~3.1 (p.~5); Eq.~33 (p.~13); Table~1 (p.~14).
\Needspace{6\baselineskip}
\section{R3-5}
\textbf{Comment.}\begin{quote}\small
Please clarify how the active-flow equations are initialized under the stated positive-equivalent-stress restriction, since the normalized deviatoric direction is undefined at the stress-free cone apex.
\end{quote}
\textbf{Response.} The revised text specifies an elastic path from the
stress-free state to the smooth yield surface. Hydrostatic elastic compression
first enters \(p>0\), \(q=0\), \(f<0\), and elastic shear then reaches the
surface at positive \(q\). The worked example gives both preload segments and
verifies positive \(q\) throughout active loading. The normalized direction is
evaluated during plastic loading, when \(q>0\). Plastic loading directly from
the apex requires a separately specified nonsmooth treatment.

\textbf{Revised locations.} Section~6 (p.~12); Eq.~28 (p.~12); Section~6.1 (p.~12).
\Needspace{6\baselineskip}
\section{R3-6}
\textbf{Comment.}\begin{quote}\small
The repeated explanations of the Flory equivalence and combined friction–dilatancy coefficient in Sections 5–8 could be shortened. A brief concluding statement of the loading and material assumptions would improve readability.
\end{quote}
\textbf{Response.} We removed the repeated algebraic summary after the
dissipation calculation and the repeated generalized-force derivation at finite
deformation. The discussion and conclusions are rewritten around the
distinction between activation, power, and closure rather than repeating the
Flory equivalence and the friction--dilatancy coefficient. The concluding
paragraph collects the assumptions: isotropy, cohesionless response, no
energetic hardening, smooth active states with \(q>0\), and monotone
compression. The worked example adds coaxial loading and its stated elastic
energy.

\textbf{Revised locations.} Section~5 (p.~8); Section~6 (p.~12); Section~7 (p.~14); Section~8 (p.~15).
\end{document}
